\section{Convoluzione}
La \textbf{convoluzione} è un'operazione importantissima all'interno del campo della teoria dei segnali; descrive come un segnale di input $f$ è
modificato da un sistema la cui risposta è $g$. Siano quindi queste due funzioni appartenenti al dominio dei complessi, la definizione è: \[
    (f*g)(t) = \int_{-\infty}^{+\infty} f(\tau)g(t-\tau) d\tau
\]

\noindent Graficamente ha luogo tenendo fermo uno dei due segnali, di norma $f$, e si inverte $g$ per poi farlo scorrere sul primo, formando un'area
in base al tempo $t$ considerato. Geometricamente, dunque, misura quanto le due funzioni si sovrappongono quando una scorre sopra l'altra. Detiene
inoltre alcune proprietà: \begin{itemize}
    \item \textbf{Commutativa}: $f*g = g*f$
    \item \textbf{Associativa}: $(f*g)*h = f*(g*h)$
    \item \textbf{Distributiva rispetto all'addizione}: $f*(g+h) = f*g + f*h$
    \item \textbf{Linearità rispetto allo scalare}: $(af)*g = a(f*g)$
\end{itemize}

\noindent Introduciamo ora il \textbf{teorema della convoluzione}, il quale la mette in relazione con la trasformata di Fourier; afferma che la
convoluzione di due segnali nel dominio del tempo è equivalente al loro prodotto nel dominio delle frequenze: \[
    \widehat{f*g} = \sqrt{2\pi}\cdot F(\omega)G(\omega)
\]

\noindent Per dimostrarlo, applichiamo la trasformata di Fourier: \begin{equation}
    \begin{split}
        \widehat{f*g}(\omega) &= \frac{1}{\sqrt{2\pi}} \int_{-\infty}^{+\infty} (f*g)(t)\cdot e^{-i\omega t} dt\\
		&= \frac{1}{\sqrt{2\pi}} \int_{-\infty}^{+\infty} \int_{-\infty}^{+\infty} f(\tau)g(t-\tau)\cdot e^{-i\omega t} d\tau dt
    \end{split}
\end{equation}

\noindent Cambiamo ora qualche variabile negli integrali; diciamo $s = t-\tau, t = st\tau$, ottenendo: \[
    e^{-i\omega t} = e^{-i\omega(s+\tau)} = e^{-i\omega s}e^{-i\omega\tau}
\]

\noindent Gli integrali ora si possono separare e per la definizione della trasformata, ognuno è uguale a $\sqrt{2\pi}$ volte la sua trasformata
normalizzata: \begin{equation}
    \begin{split}
        \widehat{f*g}(\omega) &= \frac{1}{\sqrt{2\pi}} \left(\int_{-\infty}^{+\infty} f(\tau)e^{-i\omega\tau}d\tau\right) \left(\int_{-\infty}^{+\infty}
                g(s)e^{-i\omega s}ds\right)\\
            &= \frac{1}{\sqrt{2\pi}}\cdot (\sqrt{2\pi}F(\omega)) (\sqrt{2\pi}G(\omega))\\
            &= \sqrt{2\pi}\cdot F(\omega)G(\omega)
    \end{split}
\end{equation}

\noindent A sua volta, è possibile anche applicare la convoluzione nel dominio delle frequenze, la quale corrisponderà al prodotto dei segnali nel
dominio del tempo: \[
    (F*G)(\omega) = \int_{-\infty}^{+\infty}F(\nu)G(\omega-\nu)d\nu
\]

\noindent Considerando poi l'antitrasformata di Fourier, otteniamo: \[
    \mathcal{F}^{-1}\{F*G\}(t) = \sqrt{2\pi}\cdot f(t)g(t) \implies \widehat{f\cdot g}(\omega) = \frac{1}{\sqrt{2\pi}} (F*G)(\omega)
\]

\noindent La dimostrazione inizia prendendo la trasformata del prodotto e la sostituzione della trasformata inversa per $g(t)$: \begin{equation}
    \begin{split}
        \widehat{f\cdot g}(\omega) &= \frac{1}{\sqrt{2\pi}} \int_{-\infty}^{+\infty} f(t)g(t) e^{-i\omega t} dt\\
            &= \frac{1}{\sqrt{2\pi}} \int_{-\infty}^{+\infty} f(t) \left[\int_{-\infty}^{+\infty} G(\nu)e^{i\nu t} d\nu\right]e^{-i\omega t} dt
    \end{split}
\end{equation}

\noindent Si procede invertendo l'ordine di integrazione, ottenendo un integrale di Fourier non normalizzato di $f$ nelle quadre. Siccome la formula
della trasformata include la frazione, la si aggiunge nella trasformazione, ottenendo: \begin{equation}
    \begin{split}
        \widehat{f\cdot g}(\omega) &= \frac{1}{\sqrt{2\pi}} \int_{-\infty}^{+\infty} G(\nu) \left[\int_{-\infty}^{+\infty} f(t)e^{-i(\omega-\nu) t} dt
                \right]d\nu\\
            &= \frac{1}{\sqrt{2\pi}} \int_{-\infty}^{+\infty} G(\nu)\cdot F(\omega-\nu) d\nu = \frac{1}{\sqrt{2\pi}} (F*G)(\omega)
    \end{split}
\end{equation}

%

\section{Sistemi}
Definiamo \textbf{sistema} un'entità capace di ricevere in input un segnale $x(t)$ e produrre un output $y(t)$, chiamata risposta rispetto all'input.
Ogni sistema ha al suo interno un elemento $T$ che trasforma il segnale da input ad output e può essere un dispositivo fisico, un'operazione matematica,
un filtro o molto altro ancora. Detengono alcune proprietà fondamentali:

\paragraph{Linearità:} Detta anche \textbf{superposition}, afferma che sommare o scalare l'input provoca la somma o la scalatura anche nell'output. \[
	T\{ax_i + bx_2\} \equiv aT\{x_i\} + bT\{x_2\}; \text{Con $a,b$ scalari}
\]

\noindent L'utilità di questa proprietà è fondata sul fatto che input e output possono appartenere a domini diversi; di conseguenza, se si modifica uno
dei due segnali, così farà anche l'altro senza ulteriori operazioni.

\paragraph{Tempo-invarianza:} Aggiungendo un ritardo (time-shift) al segnale di input, lo si provocherà anche in quello di output: \[
	T\{x(t-t_0)\} = y(t-t_0)
\]

\noindent Chiariamo che questa proprietà non è universale; altri sistemi possono presentare una tempo-varianza, dove in base si troverà nel tempo il
segnale di input, quello di output potrebbe essere amplificato o modificato. Per esempio $x(t-t_0) \implies tx(t-t_0)$.

\paragraph{Causalità:} Un sistema si dice causale se il suo output al tempo $t$ dipende solo dal segnale di input fino al tempo $t$ e mai dopo.
Fondamentale per poter permettere una trasformazione in tempo reale. \begin{equation}
    \begin{aligned}
        y(t) &= 2x(t) & \text{Sistema causale}\\
        y(t) &= 2x(t-1) & \text{Sistema causale}\\
        y(t) &= x(t+1) & \text{Sistema non causale}
    \end{aligned}
\end{equation}

\noindent L'ultimo sistema non è causale perché per poter riportare l'output è necessario conoscere i valori dell'input al tempo $t+1$, che è nel
futuro.

\paragraph{Memoria:} Un sistema senza memoria calcola l'output al tempo $t$ usando esclusivamente il valore dell'input al tempo $t$. Inoltre, un sistema
causale con memoria non può richiederle di conoscere il futuro, come dice la proprietà \[
	y(t) = f(x(t))
\]

\paragraph{Stabilità:} Una proprietà che vale per l'input prima della trasformazione vale anche dopo per l'output. Esistono più tipi di stabilità, come
quella asintotica, dove graficamente il segnale converge verso lo $0$ andando avanti nel tempo, oppure quella bounded-input bounded-output (BIBO), la
quale richiede che il segnale si trovi sempre all'interno di un certo intervallo.

\paragraph{Invertibilità:} Un sistema è invertibile quando il suo output contiene informazioni sufficienti per ricostruire l'input.

% TODO Sono in pari, ma il professore non ha finito di introdurre i sistemi LTI.